
\documentclass[11pt]{article}
\usepackage[margin=1in]{geometry}
\usepackage[T1]{fontenc}
\usepackage[utf8]{inputenc}
\usepackage{lmodern}
\usepackage{microtype}
\usepackage{amsmath,amssymb,mathtools,bm,amsthm}
\usepackage{hyperref}
\usepackage{booktabs,longtable,array,multirow}
\usepackage{enumitem}
\hypersetup{colorlinks=true,linkcolor=blue,citecolor=blue,urlcolor=blue}
\newtheorem{proposition}{Proposition}
\newtheorem{theorem}{Theorem}
\newtheorem{assumption}{Assumption}
\newcommand{\ii}{\mathrm{i}}
\newcommand{\dd}{\mathrm{d}}
\newcommand{\Wfour}{W_4}
\newcommand{\Qfour}{Q_4}
\newcommand{\LH}{L_H}
\newcommand{\hatG}{\widehat G}
\newcommand{\sigthree}{\sigma_3}
\newcommand{\diag}{\operatorname{diag}}
\begin{document}
\title{Technical Supplement: all-orders torus structure, phase recurrence, and numerical support}
\author{OpenAI working draft}
\date{April 18, 2026}
\maketitle

\section{Scope and status hierarchy}
This supplement accompanies the shortened main-text version of the local Type-1, $N=3$ selector reduction. It distinguishes clearly between theorem-level statements and numerically supported narrow assumptions.

\paragraph*{Theorem-level results.}
At a simple transport ramification point of the canonical exact two-sheet reduction:
\begin{enumerate}[leftmargin=2em,itemsep=2pt]
\item the exact local gauge is torus-valued to all orders;
\item the odd phase coefficients $D_{2m+1}(H)$ obey a compact Lagrange--B\"urmann recurrence;
\item the full member-dependent phase series is affine-linear in $\LH(v)$ and $\LH'(v)$ across the commuting family;
\item the same-radius local exact-WKB selector multiplier is scalar.
\end{enumerate}

\paragraph*{Numerically supported narrow assumptions used only to descend back to the original outer frame.}
\begin{assumption}[outer torus survival]
After spectator decoupling and subtraction of the universal curvature artifact, no nondecaying torus-breaking off-diagonal survives in the outer matcher.
\end{assumption}
\begin{assumption}[sector-independent residual constant part]
The residual constant non-torus factor is sector-independent and geometry-only, so it may be absorbed into the resolved-sheet convention.
\end{assumption}
The numerical basis for these assumptions is summarized in Sec.~\ref{sec:numerics}.

\section{Exact two-sheet torus theorem}
Let $v$ be a simple transport ramification point, $u_*=u(v)$, and set $u=u_*+t^2$. The reduced sheet carriers satisfy exact scalar equations
\begin{equation}
 \partial_uR_\pm=\ii(E_\pm-E_x)R_\pm,
\end{equation}
so on the $t$-cover,
\begin{equation}
 \partial_t\begin{pmatrix}R_+\\R_-\end{pmatrix}=
 \bigl(A_v(t;H)I+B_v(t;H)\sigma_3\bigr)
 \begin{pmatrix}R_+\\R_-\end{pmatrix},
\end{equation}
with
\begin{equation}
 A_v(t;H)=\ii t\,(E_++E_--2E_x),\qquad B_v(t;H)=\ii t\,(E_+-E_-).
\end{equation}
Thus the active local two-sheet system is exactly diagonal in the ordered resolved-sheet basis. Reconstruct the half-form amplitudes by
\begin{equation}
 \Psi_\pm(t;H)=\sqrt{2t}\,\Gamma_\pm(t)R_\pm(t;H)\sqrt{\dd t}.
\end{equation}
After removing the common scalar phase, the exact local basis is
\begin{equation}
 Y_v^{\mathrm{exact}}(t;H)=\diag\bigl(T_+(t)e^{\Phi_{\mathrm{rel}}(t;H)},\,T_-(t)e^{-\Phi_{\mathrm{rel}}(t;H)}\bigr),
\end{equation}
where $T_\pm(t)=\sqrt{2t}\,\Gamma_\pm(t)$ and $\Phi_{\mathrm{rel}}(t;H)=\frac{\ii}{2}S_{+-}(t;H)$. The WKB-normalized Airy basis preserves the same ordered line splitting $L_+\oplus L_-$. Consequently the matching gauge is diagonal.

\begin{theorem}
In the canonical resolved-sheet / WKB normalization, the exact local matching gauge has the form
\begin{equation}
 \hatG_v(t;H)=\exp\!\bigl(\Theta_v(t;H)\sigma_3\bigr),
\end{equation}
with analytic scalar exponent $\Theta_v(t;H)$. No torus-breaking off-diagonal term can survive in this canonical two-sheet reduction.
\end{theorem}

\begin{proof}
Both the exact and WKB-normalized local bases preserve the same ordered decomposition $L_+\oplus L_-$. Any change of basis between two ordered bases of the same direct-sum splitting must preserve each line separately, hence is diagonal. Determinant-one normalization forces the form $\exp(\Theta_v\sigma_3)$.
\end{proof}

\section{Lagrange--B\"urmann recurrence for $D_{2m+1}(H)$}
Write
\begin{equation}
 z=\lambda-v,\qquad g_v(z):=u(v+z)-u_*=z^2h_v(z),\qquad h_v(0)=\frac{u''(v)}{2}\neq 0.
\end{equation}
Choose an analytic square root and define
\begin{equation}
 \phi_v(z):=z\sqrt{h_v(z)}.
\end{equation}
Since $\phi_v'(0)\neq 0$, it has an ordinary local inverse $\psi_v$, and the two inverse branches are
\begin{equation}
 \lambda_+(t)=v+\psi_v(t),\qquad \lambda_-(t)=v+\psi_v(-t).
\end{equation}
Set
\begin{equation}
 F_H(z):=\Delta_H(v+z),\qquad \Delta_H(\lambda)=\frac{\mu(\lambda)\LH(\lambda)}{p(\lambda)},\qquad \mu^2=\Qfour.
\end{equation}
Define the odd coefficients by
\begin{equation}
 \Delta_H(\lambda_+(t))-\Delta_H(\lambda_-(t))=2\sum_{m\ge0}D_{2m+1}(H)t^{2m+1}.
\end{equation}
Ordinary Lagrange--B\"urmann inversion yields
\begin{equation}
 D_{2m+1}(H)=\frac{1}{2m+1}[z^{2m}]F_H'(z)\left(\frac{z}{\phi_v(z)}\right)^{2m+1}.
\end{equation}
Because $\phi_v(z)=z\sqrt{h_v(z)}$, this becomes
\begin{equation}
 D_{2m+1}(H)=\frac{1}{2m+1}[z^{2m}]F_H'(z)h_v(z)^{-(2m+1)/2}.
 \label{eq:kernel-supp}
\end{equation}

For computation, write
\begin{equation}
 h_v(z)=h_0\left(1+\sum_{k\ge1}a_kz^k\right),\qquad h_0=\frac{u''(v)}{2},
\end{equation}
\begin{equation}
 Q_m(z):=\left(1+\sum_{k\ge1}a_kz^k\right)^{-(2m+1)/2}=\sum_{n\ge0}q_n^{(m)}z^n.
\end{equation}
Then coefficient comparison gives
\begin{equation}
 q_0^{(m)}=1,
\end{equation}
\begin{equation}
 q_n^{(m)}=\frac{1}{n}\sum_{k=1}^n\left(\frac{1-2m}{2}k-n\right)a_kq_{n-k}^{(m)}.
 \label{eq:q-rec-supp}
\end{equation}
If $F_H'(z)=\sum_{r\ge0}\delta_r(H)z^r$, then
\begin{equation}
 D_{2m+1}(H)=\frac{h_0^{-(2m+1)/2}}{2m+1}\sum_{r=0}^{2m}\delta_r(H)q_{2m-r}^{(m)}.
 \label{eq:D-rec-supp}
\end{equation}

\section{Affine-linearity on the commuting family}
Because $\Delta_H=F\LH$ with $F(\lambda)=\mu(\lambda)/p(\lambda)$ and $\LH$ linear, we have
\begin{equation}
 \LH(v+z)=\LH(v)+\LH'(v)z.
\end{equation}
Hence
\begin{equation}
 F_H'(z)=\LH(v)F_v'(z)+\LH'(v)\bigl(F_v(z)+zF_v'(z)\bigr),\qquad F_v(z)=F(v+z).
\end{equation}
Substituting into \eqref{eq:kernel-supp} or \eqref{eq:D-rec-supp} yields
\begin{equation}
 D_{2m+1}(H)=A_m(v)\LH(v)+B_m(v)\LH'(v),
\end{equation}
with geometry-only sequences
\begin{equation}
 A_m(v)=\frac{h_0^{-(2m+1)/2}}{2m+1}\sum_{r=0}^{2m}(r+1)F_{r+1}(v)q_{2m-r}^{(m)},
\end{equation}
\begin{equation}
 B_m(v)=\frac{h_0^{-(2m+1)/2}}{2m+1}\sum_{r=0}^{2m}(r+1)F_r(v)q_{2m-r}^{(m)}.
\end{equation}
Therefore the member-dependent phase series is an explicitly computable affine-linear sequence on the full commuting family.

\section{Phase remainder and scalar multiplier}
Since $\dd u=2t\,\dd t$, the exact phase mismatch is
\begin{equation}
 S_{+-}(t;H)=4\sum_{m\ge0}\frac{D_{2m+1}(H)}{2m+3}t^{2m+3}.
\end{equation}
The linear Langer coordinate removes the cubic term exactly, so
\begin{equation}
 \Theta_v^{\mathrm{ph}}(t;H)=2\ii\sum_{m\ge1}\frac{D_{2m+1}(H)}{2m+3}t^{2m+3}.
\end{equation}
If
\begin{equation}
 \Theta_v^{\mathrm{ph}}(t;H)=\sum_{n\ge2}\Xi_{2n+1}(H)t^{2n+1},
\end{equation}
then
\begin{equation}
 \Xi_{2n+1}(H)=\frac{2\ii}{2n+1}D_{2n-1}(H),\qquad n\ge2.
\end{equation}
Thus
\begin{equation}
 \Xi_5(H)=\frac{2\ii}{5}D_3(H),\qquad \Xi_7(H)=\frac{2\ii}{7}D_5(H),\qquad \Xi_9(H)=\frac{2\ii}{9}D_7(H),\dots
\end{equation}
Because the local inverse branches and compositions are analytic up to the nearest singularity, the local series is convergent there; the local Borel sum therefore coincides with the analytic function itself. In the WKB-normalized Airy convention, the active same-radius $S_1\to S_0$ crossing carries the scalar multiplier
\begin{equation}
 \mathfrak s_v^{\mathrm{BS}}(t;H)=-\ii\exp\!\bigl(2\Theta_v(t;H)\bigr).
\end{equation}

\section{Canonical $v_4$ example and explicit $t^7$ result}
For the standard real family
\begin{equation}
 (\epsilon_0,\epsilon_1,\epsilon_2)=(-2,0,3),\qquad (\gamma_0,\gamma_1,\gamma_2)=\left(1,\frac45,\frac65\right),
\end{equation}
at the canonical upper-half-plane selector point $v_4=-0.8019106146-1.1076910061\,\ii$, the recurrence reproduces the known $\Xi_5$ and yields explicit higher coefficients. In affine-linear form,
\begin{equation}
 \Xi_7(H)=(-0.1218869377+0.5875264733\,\ii)\,\beta_H+(-0.5377433968+1.2623260628\,\ii)\,\alpha_H,
\end{equation}
\begin{equation}
 \Xi_9(H)=(-1.0157641727+0.8877006452\,\ii)\,\beta_H+(-2.5833648230+1.4527859580\,\ii)\,\alpha_H.
\end{equation}
For the resonant wall member used in the local selector notes, this gives
\begin{equation}
 \Xi_5(H_{\mathrm{wall}})=0.0511662458609726+0.106488885281404\,\ii,
\end{equation}
\begin{equation}
 \Xi_7(H_{\mathrm{wall}})=-0.0590453140023154+0.151640061843118\,\ii,
\end{equation}
\begin{equation}
 \Xi_9(H_{\mathrm{wall}})=-0.302262984767034+0.183667018646898\,\ii.
\end{equation}

\section{Numerical analyses that led to the outer canonical form}
\label{sec:numerics}
Table~\ref{tab:numerics} records the numerical path that led to the outer canonical form. The theorem-level statements of the main text do not rely on these data; the data are used only to justify the narrow descent assumptions that identify the canonical two-sheet gauge with the original outer matcher.

\begin{table}[h]
\caption{High-precision assays used to justify the outer canonical form.}
\label{tab:numerics}

\begin{tabular}{p{0.28\linewidth}p{0.18\linewidth}p{0.39\linewidth}}
Assay & Accepted samples & Outcome / interpretation\\
\hline
Resolved-sheet $\hatG_v$ off-diagonal probe & 1200 & No off-diagonal counterexample found: 1200/1200 ``no''. Supports the exact two-sheet torus theorem from the numerical side.\\
Spectator-decoupled torus survival & 1001 (301 targeted high-curvature) & 1001/1001 pass at the main threshold. All 9/9 legacy high-curvature failures from the earlier raw assay pass on retest. Supports Assumption 1.\\
Sector-independence of residual constant part & 1001 (301 targeted high-curvature) & 999/1001 base-run passes; only 2 borderline extreme-curvature misses. Both pass refined 100-digit retest, leaving no durable counterexample. Supports Assumption 2.\\
Lagrange-inversion recurrence validation & 1200 & 1200/1200 pass at $10^{-2}$ relative threshold and 1198/1200 at $10^{-3}$ for $(D_1,D_3,D_5)$ simultaneously.\\
\end{tabular}

\end{table}

Two details matter. First, the spectator-decoupled torus-survival assay intentionally oversampled the extreme-curvature regime, the most plausible source of counterexamples. Second, the two base-run sector-independence misses were not stable under refinement. Their refined retest values were:
\begin{center}
\begin{tabular}{rrrrr}
\toprule
sample & $|u''(v)|$ & $|c_0|_{\rm refined}$ & fit error & pass at $10^{-6}$?\\
\midrule
810 & 370.801 & 7.14685e-07 & 3.53439e-07 & yes\\
830 & 299.187 & 6.42155e-07 & 3.26301e-07 & yes\\
\bottomrule
\end{tabular}
\end{center}
These data strongly support the view that the surviving local structure in the original frame is torus-valued up to a sector-independent geometry-only constant factor.

\section{What remains open and high-leverage tactics}
The remaining unknowns are now sharply concentrated.
\begin{enumerate}[leftmargin=2em,itemsep=2pt]
\item \textbf{Selector sufficiency.} The necessary sieve---transport ramification plus phase resonance---is formalized, but one still needs a proof that every simple phase-active transport ramification point actually inserts the scalar local block.
\item \textbf{Exact insertion rule.} The scalar local block is now explicit in the canonical two-sheet normalization. What remains is the exact local-to-global comparison theorem that inserts it into the continuation functor.
\item \textbf{Transport recurrence.} The phase side now has a closed affine-linear recurrence. A matching all-orders recurrence for the transport exponent $\Theta_v^{\mathrm{tr}}(t)$ would place the full torus exponent on the same algorithmic footing.
\item \textbf{Nongeneric collisions.} Degenerate transport ramification and common zeros of $\Qfour$ and $\Wfour$ still require a separate analysis.
\end{enumerate}

The highest-leverage next moves are symbolic: derive the transport recurrence, prove selector sufficiency in the scalarized two-sheet model, and then prove the exact global insertion rule. Numerics should now be used only as theorem-shaping tools: stress-test sufficiency hypotheses in the extreme-curvature regime, guess closed recurrences, and isolate the precise nongeneric hypotheses needed for theorem statements.

\begin{thebibliography}{99}
\bibitem{summary} \emph{Multistate Landau--Zener Solver: Augmented ODE with Gauged Dynamics} (internal summary).
\bibitem{gap} \emph{A note on the sextic gap polynomial for $N=3$ Type-1 matrices} (internal note).
\bibitem{primer} \emph{Type-1, $N=3$ Landau--Zener primer: double cover, reduced ODE, gluing theorem, and selector} (internal primer).
\bibitem{strategy} \emph{Formalized strategy diagram for Type-1, $N=3$ Landau--Zener} (internal note).
\bibitem{match} \emph{Local selector matching at a transport ramification point in Type-1, $N=3$} (internal note).
\bibitem{gauge} \emph{Transport versus phase in the local gauge expansion at a simple transport ramification point} (internal note).
\end{thebibliography}
\end{document}
